% =====================================================================
% 附录 A：理论。由原 appendix.tex（A.1 失败诊断 / A.2 证明 / A.5 附加理论结果）
% 与 appendix_additional.tex（B.1 延期理论陈述，前移供证明引用）重组而来。
% =====================================================================
\section{Proofs and theoretical analysis}
\label{app:theory}

\subsection{Proofs of the main guarantees}

\begin{lemma}[OPB KL decomposition]
\label{lem:opb-kl-decomposition}
For a diagonal posterior covariance and the prior of Eq.~\eqref{eq:opb-prior}, the
standard diagonal-Gaussian KL formula reads
\begin{equation}
\KL\bigl(q_\phi(z|x)\,\|\,p_\theta^{\mathrm{OPB}}(z|y)\bigr)
= \frac{1}{2\tau^2}\|\mu_\phi(x) - a\,q_{\theta,y}\|^2
+ \frac12 \sum_{r=1}^{d}\Bigl[
\frac{\sigma_{\phi,r}^2(x)}{\tau^2} - 1 - \log\frac{\sigma_{\phi,r}^2(x)}{\tau^2}
\Bigr],
\label{eq:opb-kl-decomposition}
\end{equation}
a class-anchor mismatch term plus a variance mismatch term that is non-negative and
vanishes exactly at $\sigma_{\phi,r}^2(x) = \tau^2$ for all $r$.
\end{lemma}

\begin{proof}[Proof of Theorem~\ref{prop:pgib-alignment}]
The Gaussian KL of Lemma~\ref{lem:opb-kl-decomposition} is the sum of its mean and
variance terms, the latter non-negative, so with the fixed prior variance $\tau^2=1$,
\begin{equation*}
\KL\bigl(q_\phi(z|x)\,\|\,p_\theta(z|y)\bigr)
\ge\frac1{2\tau^2}\|\mu_\phi(x)-m_\theta(y)\|^2
=\frac{\|\mu_\phi(x)-m_\theta(y)\|^2}{2},
\end{equation*}
and hence $\mathcal D(\widehat w)\le\mathcal K(\widehat w)$. For any feasible $w_0$
in the same family, $\varepsilon$-optimality of $\widehat w$, $\ell\ge0$, and
$\beta>0$ give
\begin{equation*}
\beta\mathcal K(\widehat w)
\le \mathcal R(\widehat w)+\beta\mathcal K(\widehat w)
\le \mathcal R(w_0)+\beta\mathcal K(w_0)+\varepsilon.
\end{equation*}
Division by $\beta$ proves Eq.~\eqref{eq:pgib-comparator-bound}.
\end{proof}

\begin{proof}[Proof of Theorem~\ref{thm:pgib-separation}]
Jensen's inequality gives
\begin{equation*}
\|m(y) - m_\theta(y)\| \le \sqrt{2\tau^2 D(y)}
\end{equation*}
(the norm is convex, and $\E\|\mu_\phi(X)\|^2<\infty$ makes $m(y)$ well defined),
and the triangle inequality on
\begin{equation*}
m(y) - m(y') = (m(y) - m_\theta(y)) + (m_\theta(y) - m_\theta(y'))
+ (m_\theta(y') - m(y'))
\end{equation*}
gives the two-level bound of
Eq.~\eqref{eq:pgib-separation}, the last step
by the definition of $\Delta(y,y')$. The first level is positive exactly
when $\Gamma_{yy'}<1$ and the declared-margin level exactly when
$\sqrt{2\tau^2D(y)}+\sqrt{2\tau^2D(y')}<\Delta(y,y')$; for OPB the
anchor pair distance is the constant $\sqrt2a$ of
Eq.~\eqref{eq:opb-prior}, so the declared bound satisfies $\Delta(y,y')=\sqrt2a$ and
the two conditions coincide. With equal mismatch $D$ the common condition
\begin{equation*}
\Gamma_{yy'}=\frac{2\sqrt{2\tau^2D}}{\sqrt2a}<1
\qquad\Longleftrightarrow\qquad
D<\frac{a^2}{4\tau^2}.
\end{equation*}
For continuous labels the Jensen inequality holds
pointwise for $P_Y$-a.e.\ $y$ (regular conditional expectation), and intersecting
the two measure-one sets of valid $y$'s yields the $P_Y\otimes P_Y$-almost-everywhere
pair statement.
\end{proof}

\begin{remark}[Finite-sample estimation of the continuous case]
\label{rem:pgib-binning}
On finite samples the regular conditional expectations of the continuous case are
estimated by
binning the label axis into intervals $B_i$ with representative labels $\bar y_i$
and centers $m_i = \E[\mu_\phi(X) \mid Y \in B_i]$; the bin diameter adds a
quantization error bounded by
\begin{equation*}
\rho_{\max}\,\operatorname{diam}(B_i) + \rho_{\max}\,\operatorname{diam}(B_j)
\end{equation*}
(by the upper bound of Eq.~\eqref{eq:pgib-lipschitz}), which vanishes as the bins
shrink. The separation statement itself is pointwise and needs no binning.
\end{remark}

\begin{proof}[Proof of Corollary~\ref{cor:pgib-beta-transfer}]
Theorem~\ref{prop:pgib-alignment} gives
\begin{equation*}
\mathcal D(\widehat w_\beta)\le B_\beta.
\end{equation*}
For discrete labels, the tower property gives
$\mathcal D(\widehat w_\beta)=\sum_k\pi_kD(k)$, hence
\begin{equation*}
D(k)\le B_\beta/\pi_k\le B_\beta/\pi_{\min};
\end{equation*}
for continuous labels, Markov's inequality on the nonnegative random variable
$D(Y)$ gives
\begin{equation*}
P_Y\bigl(D(Y)>B_\beta/\xi\bigr)\le\xi.
\end{equation*}
In both cases
Theorem~\ref{thm:pgib-separation} then gives
Eq.~\eqref{eq:pgib-beta-separation}. The rate statement follows
directly from the definition of $B_\beta$: with $\mathcal K(w_0)=0$ and
$\sup_\beta\varepsilon_\beta<\infty$, $B_\beta=O(\beta^{-1})$ so the certified
deficit is $O(\beta^{-1/2})$; with $\mathcal K(w_0)>0$,
$B_\beta\to\mathcal K(w_0)$ as $\beta\to\infty$.
\end{proof}

\begin{remark}[What the framework constrains (and what it does not)]
\label{rem:counterexample}
The constraints of the geometric prior bind the reference, not the deployed
representation: for any orthonormal frame $\{q'_k\}$ and the class-constant posterior
$q(z|Y=k) = \mathcal{N}(a\,q'_k, \tau^2 I_d)$, the trainable frame can rotate
to match ($q_{\theta,k} \to q'_k$), yielding $\KL = 0$, $I(X;Z|Y) = 0$, and
$I(Y;Z) > 0$: a full-retention label code with a zero certificate. The framework
therefore constrains the geometric \emph{form} of the label code (precisely what CEB
left unspecified), and does not claim that the certificate meters retained label
information (\citealp{kolchinsky2019caveats}). In the GPB instance used by the real-data protocol, the remaining
freedom at certificate zero is the frame rotation only.
\end{remark}

\emph{Scope of the comparison argument.} Theorem~\ref{prop:pgib-alignment} applies to
arbitrary $P_{XY}$ and every feasible comparator in the fitted model class; it does not
require $q(z|x)=p_\theta(z|y)$. Its usefulness depends on the comparator KL. Label
ambiguity, limited encoder capacity, and covariance misspecification produce the
non-zero floor made explicit in Proposition~\ref{prop:pgib-ambiguity-decomposition}.
The stronger statement that this floor is zero still requires the assumptions of
Corollary~\ref{cor:pgib-exact-alignment}: deterministic labels, exact posterior--prior
matching, and a compatible head. We do not assume those exact-realizability premises
for the real-data benchmarks. Their performance comparison evaluates utility, while
their geometric diagnostics report realized mismatch and geometry. The controlled
known-noise experiment instead evaluates the ambiguity term where its population value
is analytic, including the deterministic distribution as its zero-floor special case.

\subsection{Deferred theory statements and proofs}
\label{app:deferred-theory}

\begin{proposition}[Ambiguity--approximation decomposition]
\label{prop:pgib-ambiguity-decomposition}
For OPB classification with $\eta_k(x)=P(Y=k\mid X=x)$ and a comparator mean $h(x)$,
\begin{equation}
\mathcal K(w_0)=\frac{a^2}{2\tau^2}\E[1-\|\eta(X)\|_2^2]
+\frac1{2\tau^2}\E\left\|h(X)-a\sum_k\eta_k(X)q_{0,k}\right\|^2+V_0,
\label{eq:opb-ambiguity-decomposition}
\end{equation}
where $V_0$ is the covariance mismatch. For EPB regression,
\begin{equation}
\mathcal K(w_0)=\frac{\rho^2}{2\tau^2}\E[\operatorname{Var}(\tilde Y\mid X)]
+\frac1{2\tau^2}\E\|h(X)-\rho\nu(X)u_0\|^2+V_0.
\label{eq:epb-ambiguity-decomposition}
\end{equation}
\end{proposition}

\begin{proof}
For OPB let $A_Y:=a q_{0,Y}$. Conditional on $X$, orthonormality gives
$\E[A_Y\mid X]=a\sum_k\eta_k(X)q_{0,k}$ and
\begin{equation*}
\E\!\left[\|A_Y-\E[A_Y\mid X]\|^2\mid X\right]
=a^2\left(1-\|\eta(X)\|_2^2\right).
\end{equation*}
The conditional bias--variance identity then yields
\begin{equation*}
\E\|h(X)-A_Y\|^2
=a^2\E[1-\|\eta(X)\|_2^2]
+\E\left\|h(X)-a\sum_k\eta_k(X)q_{0,k}\right\|^2.
\end{equation*}
Adding the diagonal-Gaussian covariance term $V_0$ of
Lemma~\ref{lem:opb-kl-decomposition} proves
Eq.~\eqref{eq:opb-ambiguity-decomposition}. For EPB set $A_Y:=\rho\tilde Y u_0$; then
\begin{equation*}
\E[A_Y\mid X]=\rho\nu(X)u_0
\qquad\text{and}\qquad
\E\bigl[\|A_Y-\E[A_Y\mid X]\|^2\mid X\bigr]
=\rho^2\operatorname{Var}(\tilde Y\mid X).
\end{equation*}
The same identity plus $V_0$ proves Eq.~\eqref{eq:epb-ambiguity-decomposition}.
\end{proof}

Substituting either decomposition into $B_\beta$ in
Corollary~\ref{cor:pgib-beta-transfer} separates the comparator-dependent geometry
transfer floor into label ambiguity, encoder mean approximation, and covariance
mismatch.

\begin{corollary}[Zero-floor exact-realizability special cases]
\label{cor:pgib-exact-alignment}
Under deterministic labels, exact posterior--prior matching, and a compatible Bayes
head, $\mathcal K(w_0)=0$ and Theorem~\ref{prop:pgib-alignment} gives
\begin{equation}
\mathcal D(\widehat w)\le(\ln K+\varepsilon)/\beta
\label{eq:pgib-alignment}
\end{equation}
for classification. For EPB regression,
\begin{equation}
\mathcal D(\widehat w)\le(\tau^2/\rho^2+\varepsilon)/\beta.
\label{eq:pgib-alignment-regression}
\end{equation}
\end{corollary}

\begin{proof}
For classification, the assumed comparator $q(z|x)=p_\theta(z|y)$ has
$\mathcal K(w_0)=0$. Its Bayes cross-entropy is
\begin{equation*}
C_{\mathrm G}=H(Y\mid Z_{\mathrm{align}})\le\ln K,
\end{equation*}
since with shared covariance the determinant factors cancel and its posterior is the
anchor-energy rule. Applying Eq.~\eqref{eq:pgib-comparator-bound} gives
Eq.~\eqref{eq:pgib-alignment}. For regression the aligned comparator again has zero
KL. Under the tied head its error is $(\tau/\rho)\eta$ with
$\eta\sim\mathcal N(0,1)$, whose standardized squared risk is
\begin{equation*}
\tau^2/\rho^2;
\end{equation*}
the same proposition gives Eq.~\eqref{eq:pgib-alignment-regression}.
\end{proof}

\begin{proposition}[Positive rate cost of class-independent means]
\label{prop:opb-collapse}
If $\E[\mu_\phi(X)\mid Y=k]=c$ for every class, then
\begin{equation}
\E\left[\frac{\|\mu_\phi(X)-a q_{\theta,Y}\|^2}{2\tau^2}\right]
\ge\frac{a^2}{2\tau^2}\left(1-\sum_k\pi_k^2\right).
\label{eq:opb-collapse}
\end{equation}
\end{proposition}

\begin{proof}
By Jensen's inequality the left side is at least
\begin{equation*}
\tfrac{1}{2\tau^2}\sum_k \pi_k \|c - a q_{\theta,k}\|^2
= \tfrac{1}{2\tau^2}\Bigl(\|c\|^2 - 2a\,c^{\top}\!\sum_k\pi_k q_{\theta,k} + a^2\Bigr).
\end{equation*}
Minimizing over $c$ gives
\begin{equation*}
c = a\sum_k\pi_k q_{\theta,k},
\end{equation*}
and by orthonormality $\|\sum_k\pi_k q_{\theta,k}\|^2 = \sum_k \pi_k^2$, yielding
Eq.~\eqref{eq:opb-collapse}.
\end{proof}

\begin{corollary}[The orthogonal instance instantiates conditional posterior separation]
\label{prop:opb-separation}
For OPB, $\tau^2=1$ and $\Delta(y,y')=\sqrt2a$, so
\begin{equation}
\|m(y)-m(y')\|\ge\sqrt2a-\sqrt{2\tau^2D(y)}-\sqrt{2\tau^2D(y')}.
\label{eq:opb-separation}
\end{equation}
\end{corollary}
\subsection{Additional theoretical results}

\begin{corollary}[Label-information lower bound in the aligned Gaussian model]
\label{cor:opb-fano}
Assume uniform classes and exact alignment
$q(z|Y=k) = \mathcal{N}(a\,q_{\theta,k}, \tau^2 I_d)$. The pairwise error probability
of the nearest-anchor classifier is
\begin{equation}
P_{\mathrm{pair}} = \Phi\Bigl(-\frac{a}{\sqrt{2}\,\tau}\Bigr),
\qquad
P_e \le \min\bigl\{1, (K-1)P_{\mathrm{pair}}\bigr\},
\end{equation}
and whenever the displayed bound satisfies $\bar P_e \le 1 - 1/K$, Fano's inequality
gives
\begin{equation}
I(Y;Z) \ge \log K - h_2(\bar P_e) - \bar P_e \log(K-1).
\end{equation}
This links the anchor scale-to-noise ratio $a/\tau$ to the retained label
information, conditionally on the aligned Gaussian model.
\end{corollary}

\begin{proof}
Two anchors have distance $\sqrt{2}a$ by Proposition~\ref{prop:opb-frame}; the
equal-covariance binary decision boundary is halfway between them, giving
$\Phi(-\sqrt{2}a/(2\tau))$. The union bound and Fano's inequality on
$[0, 1-1/K]$ conclude.
\end{proof}

In the deterministic case ($Y=f(X)$), the CEB objective is minimized only at the
copy corner $Z=Y$ for every weight $\gamma>0$, so no interior point of the IB curve
is ever a Lagrangian minimizer (\citealp{kolchinsky2019caveats}, Eq.~7). The Caveats
paper repairs this by squaring the compression term; the anchored reference supplies
the same strict convexity by a different mechanism, without squaring, through the
quadratic displacement term built into the KL.

\begin{proposition}[Joint strict convexity and unique optimum on a restricted aligned surrogate]
\label{prop:opb-sweep}
Consider the class-symmetric aligned family: every input of class $y$ is mapped to the
posterior $q_y = \mathcal{N}(s\,q_{\theta,y}, \sigma^2 I)$ ($s \ge 0$, $\sigma^2 > 0$;
anchor scale $a > 0$), and the classifier is the anchor-energy classifier of
Proposition~\ref{prop:opb-energy} (weights $w_j = c\,q_{\theta,j}$ with scale
$c = a/\tau^2$ fixed by construction). With the \emph{stochastic training code}
$z \sim q_y$ (as in the actual objective, Eq.~\eqref{eq:pgib-objective}), the
training cross-entropy is
\begin{equation}
g(s, \sigma^2) := \E_\eta\Bigl[
\ln\Bigl(1 + \textstyle\sum_{j \ne y} e^{-cs + c\sigma(\eta_j - \eta_y)}\Bigr)
\Bigr],
\qquad \eta \sim \mathcal{N}(0, I_K),
\label{eq:opb-ce-train}
\end{equation}
and the training Lagrangian, with the KL term of
Lemma~\ref{lem:opb-kl-decomposition}, is
\begin{equation}
L_\beta(s, \sigma^2) = g(s, \sigma^2)
+ \beta\Bigl[
\tfrac{1}{2\tau^2}(s-a)^2
+ \tfrac{d}{2}\bigl(\sigma^2/\tau^2 - 1 - \ln(\sigma^2/\tau^2)\bigr)
\Bigr].
\label{eq:opb-lag-train}
\end{equation}
Then for every $\beta > 0$, $L_\beta$ is strictly convex jointly in $(s, \sigma)$ and
has a unique joint minimizer $(s^*(\beta), \sigma^*(\beta))$, interior in $s$ since
$a > 0$: the stochastic training objective has a well-defined optimum, in contrast to
the deterministic IB Lagrangian, whose optimum is pinned at a corner for every $\beta$
(\citealp{kolchinsky2019caveats}). This is a \emph{restricted surrogate result}: it
analyzes two scalar variables on the class-symmetric aligned family, with the frame,
the anchor scale, and the classifier scale held fixed; the full OPB objective, with
its nonlinear posterior, trainable frame, and non-convex degrees of freedom, lies
outside this analysis, and the result neither establishes $\beta$-tunability of the
full objective nor makes any claim on the information plane.
The deterministic evaluation code ($z = \mu$) is the $\sigma \to 0$ limit, and
$g \ge \ln(1 + (K-1)e^{-cs})$ by Jensen (the log-sum-exp is convex), so the statement
strictly generalizes the deterministic surrogate analysis. We do \emph{not} claim here
that the map $\beta \mapsto s^*(\beta)$ is monotone or that the sweep is one-to-one:
the minimizer $\sigma^*$ responds to $\beta$ through the cross-derivative
$\partial^2 g/\partial s\,\partial\sigma$, and establishing the sign of
$\mathrm{d}s^*/\mathrm{d}\beta$ requires the full two-dimensional stationarity
analysis, which we leave open; the $\beta$-sweep is reported empirically in the
experiments. The fixed-scale condition holds for the main model by construction
(Proposition~\ref{prop:opb-energy}), so the statement describes the actual training
objective up to the aligned-family idealization. We do not analyze the free-classifier
variant on this surrogate: under the stochastic training code, a diverging classifier
scale does not produce a degenerate solution: at $\operatorname{KL} \to 0$ the
posterior variance is pinned at $\sigma^2 = \tau^2 > 0$, so the class-conditional
Gaussians overlap and the training cross-entropy of a free classifier diverges with
its scale rather than vanishing (for $K = 2$ it contains
$\E[\ln(1 + e^{-c(a + \tau\xi)})]$, whose integrand grows linearly in $c$ on the
event $\xi < -a/\tau$, which has positive probability). An analysis of the free
classifier would require an explicit classifier-norm constraint or regularizer and is
left to future work.
\end{proposition}

\begin{proof}
For each fixed $\eta$, the integrand of Eq.~\eqref{eq:opb-ce-train} is a log-sum-exp of
functions affine in $(s, \sigma)$ (the constant $1$ has slope zero and each exponential
has slope $(-c,\, c(\eta_j - \eta_y))$), hence jointly convex in $(s, \sigma)$, and the
expectation $g$ is jointly convex. No strictness is claimed for $g$: for $K = 2$ the
integrand depends on $(s, \sigma)$ only through the single linear combination
$-cs + c\sigma(\eta_2 - \eta_1)$, so its Hessian has rank one. Write the KL term as
\begin{equation*}
R(s, \sigma) := \tfrac{1}{2\tau^2}(s-a)^2
+ \tfrac{d}{2}\bigl(\sigma^2/\tau^2 - 1 - \ln(\sigma^2/\tau^2)\bigr);
\end{equation*}
its Hessian in $(s, \sigma)$ is
\begin{equation*}
\nabla^2 R = \operatorname{diag}\Bigl(\tau^{-2},\; d\,\bigl(\tau^{-2} + \sigma^{-2}\bigr)\Bigr),
\end{equation*}
strictly positive definite for every $\sigma > 0$, so $\beta R$ is strictly convex in
$(s, \sigma)$. Hence $L_\beta = g + \beta R$ is the sum of a jointly convex and a
strictly convex function, and is strictly convex in $(s, \sigma)$. Since
$R \to +\infty$ as $s \to \infty$, $\sigma \to 0^+$, or $\sigma \to \infty$, the
sublevel sets of $L_\beta$ are compact, and the unique joint minimizer exists (by
strict convexity it is unique). For $a > 0$ the minimizer is interior in $s$: at
$s = 0$,
\begin{equation*}
\frac{\partial L_\beta}{\partial s}\Big|_{s=0}
= -c\,\E[p_{\mathrm{w}}] - \frac{\beta a}{\tau^2} < 0,
\end{equation*}
where $p_{\mathrm{w}}$ is the error probability of the sampled code, so $s^* > 0$.
\end{proof}

\begin{remark}[The sweep is a KL--CE surrogate curve; scope and limitations]
\label{rem:opb-scope}
Proposition~\ref{prop:opb-sweep} proves joint strict convexity and uniqueness of the
optimum for the stochastic training objective in the KL--CE plane on the aligned family
at fixed classifier scale: a restricted surrogate result, which does not establish
$\beta$-tunability of the full OPB objective, and in any case not a statement on the
information plane of the Caveats paper; the results of the main text
establish a structured conditional prior, an exact KL geometry, and conditional
separation guarantees, and they do \emph{not} establish a mutual information upper
bound tighter than CEB's: Eq.~\eqref{eq:gap} is the usual CEB bound, here
for a particular prior family. The orthogonal-frame construction fixes all pairwise anchor
distances, so the framework should not claim to recover arbitrary semantic class
similarities. All geometric statements are conditional on
Assumption~\ref{ass:fullrank}: orthonormality holds exactly by the polar construction
wherever the factorization exists, and the condition is verified, not imposed,
along the reported runs (Section~\ref{sec:opb}); the construction is exactly
equivariant under relabeling of the classes by Eq.~\eqref{eq:opb-equivariance}, so
no class index plays a distinguished role. If the frame were computed from the
current input batch rather than from the all-class prior matrix, the prior would
become batch-dependent and the certificate-preservation argument would no longer
follow directly. Finally, the
anchor-energy classifier of Proposition~\ref{prop:opb-energy} fixes the classifier
scale to $a/\tau^2$ by construction, removing the independent classifier-margin
escape route in the aligned surrogate of Proposition~\ref{prop:opb-sweep}; whether a
free classifier (with an explicit norm constraint or regularizer) admits its own
degenerate route is left to future work, and the free linear classifier is kept in
the experiments as a variant for comparison.
Whether the joint optimum of Proposition~\ref{prop:opb-sweep} moves monotonically in
$\beta$ is likewise open: it requires the full two-dimensional stationarity analysis,
in which the optimal variance responds to $\beta$ through the cross-derivative
$\partial^2 g/\partial s\,\partial\sigma$ of the stochastic cross-entropy.
The trainable frame and axis constitute a rotation channel that this paper does not
address: the certificate's zero set contains the full-retention orthogonal label codes
at any frame rotation (the reference can rotate to meet them), and the framework pins
the \emph{geometry of the code}; it does not claim that the certificate reads off
the retained label information. The trainability of the frame and axis is examined
empirically in the frozen-prior ablation of Appendix~\ref{app:ablation}
(Tables~\ref{tab:ablation-ci} and~\ref{tab:ablation-ci-epb}).
\end{remark}
